\section{CVP Comparison and Kernel Ambiguity}

CVP asks for a closest vector in an explicitly presented lattice
\cite{micciancioGoldwasser2002}. HCP asks for
the projection of \(b\) into the kernel of the operative hidden Hodge operator.

\subsection{Known-Lattice CVP}

An explicit Euclidean CVP instance consists of a full-column-rank generator
matrix \(B\in\mathbb R^{m\times n}\) and a target \(t\in\mathbb R^m\). Its
discrete lattice is
\[
  \Lambda_B=B\mathbb Z^n=\{Bz:z\in\mathbb Z^n\}.
\]
The injectivity of \(z\mapsto Bz\) records that the displayed columns are a
lattice basis. A correct closest-vector answer is a vector
\(v\in\Lambda_B\) such that
\[
  \|t-v\|^2\le \|t-w\|^2
  \qquad
  \text{for every }w\in\Lambda_B.
\]
Thus, for fixed dimensions, a CVP oracle \(O\) is correct when
\[
\begin{aligned}
  \mathsf{CVP}_{\mathrm{corr}}(O)
  &\quad\Longleftrightarrow\quad
  \forall\text{ full-column-rank }B,t,\\
  &\qquad O(B,t)\in\Lambda_B
  \quad\text{and}\quad
  \forall w\in\Lambda_B,\\
  &\qquad \|t-O(B,t)\|^2\le \|t-w\|^2.
\end{aligned}
\]
The weaker output condition is merely lattice membership:
\[
  \forall B,t,\ O(B,t)\in\Lambda_B .
\]
Fix a class \(\mathcal C\) of admissible algorithms, such as polynomial-time or
probabilistic polynomial-time algorithms. The external CVP hardness assumption
is the schema
\[
\begin{aligned}
  \mathsf{CVP}_{\mathrm{hard}}(\mathcal C)
  \quad\Longleftrightarrow\quad
  \neg\exists O,\ 
    O\in\mathcal C\wedge\mathsf{CVP}_{\mathrm{corr}}(O).
\end{aligned}
\]

\subsection{Presented Linear Images}

Let \(R\) be a field. The field-linear image of \(B\) is a different object:
\[
  \mathcal V_B(R)=BR^n=\{Bx:x\in R^n\}.
\]
Over a field this is a linear subspace, not the discrete lattice
\(B\mathbb Z^n\).

Let \(B\) present \(\mathcal V_B(R)\). If a matrix
\(\Delta\) satisfies \(\ker(\Delta)=\mathcal V_B(R)\), then any solver satisfying
\[
  \Delta S(\Delta,t)=0
\]
returns \(S(\Delta,t)\in\mathcal V_B(R)\).

\begin{proposition}[Kernel realization for explicit generated subspaces]\label{prop:hcp-to-cvp}
  Let \(B\in R^{n\times n}\) be an explicit square generator matrix over a
  field and set
  \(\mathcal V_B(R)=\{Bx:x\in R^n\}\). There exists a matrix \(\Delta\) such that
  \[
    \ker(\Delta)=\mathcal V_B(R) .
  \]
  Consequently, any solver \(S(\Delta,b)\) whose output is
  \(\Delta\)-harmonic for every
  presented \(\Delta\) returns a vector in the explicit generated subspace
  \(\mathcal V_B(R)\) when run on this realized operator.
\end{proposition}

\begin{proof}
  Let \(B\in R^{n\times n}\) be an explicit generator matrix over a field
  \(R\), set
    \[
      \mathcal V_B(R)=\{Bx:x\in R^n\}\subseteq R^n,
    \]
    and let \(t\in R^n\) be the solver input. We construct a matrix \(\Delta\)
    with
  \[
    \ker(\Delta)=\mathcal V_B(R).
  \]
  Put \(V=\mathcal V_B(R)\). Since \(R^n\) is finite-dimensional over a field,
  choose a complementary subspace \(W\) with
  \[
    V+W=R^n,
    \qquad
    V\cap W=0 .
  \]
  Thus every \(y\in R^n\) has a unique decomposition
  \[
    y=v+w,
    \qquad
    v\in V,\quad w\in W.
  \]
  Equivalently, there is a unique element \(w_y\in W\) such that
  \(y-w_y\in V\). Define
  \[
    P(y)=w_y.
  \]
  The uniqueness of the \(V\oplus W\) decomposition gives linearity: for
  \(y,z\in R^n\), the vector \(P(y)+P(z)\) lies in \(W\), and
  \[
    (y+z)-(P(y)+P(z))
    =
    (y-P(y))+(z-P(z))\in V,
  \]
  so \(P(y+z)=P(y)+P(z)\). Similarly, for \(a\in R\),
  \(aP(y)\in W\) and
  \[
    ay-aP(y)=a(y-P(y))\in V,
  \]
  so \(P(ay)=aP(y)\).

  We now compute the kernel. If \(y\in V\), then \(0\in W\) and \(y-0\in V\),
  hence uniqueness gives \(P(y)=0\). Conversely, if \(P(y)=0\), then
  \(y-P(y)=y\in V\). Therefore
  \[
    \ker P=V=\mathcal V_B(R).
  \]
  Let \(\Delta\) be the matrix of the linear map \(P\). Then, for every
  \(y\in R^n\),
  \[
    \Delta y=0
    \quad\Longleftrightarrow\quad
    y\in\ker P
    \quad\Longleftrightarrow\quad
    y\in\mathcal V_B(R)
    \quad\Longleftrightarrow\quad
    \exists x,\ Bx=y .
  \]

    Let \(S:R^{n\times n}\times R^n\to R^n\) be a presented-operator Hodge
  solver satisfying
  \[
    \forall \Delta\, b,\qquad
    \Delta(S(\Delta,b))=0.
  \]
    Applying this to the \(\Delta\) constructed above and \(b=t\) gives
  \[
    \Delta(S(\Delta,t))=0.
  \]
  By the kernel calculation, there exists \(x\) with
  \[
    Bx=S(\Delta,t).
  \]
    Hence \(S(\Delta,t)\in\mathcal V_B(R)\).
\end{proof}

When \(R=\mathbb R\), if a presented-operator solver additionally minimizes
Euclidean distance among \(\Delta\)-harmonic vectors, the same kernel equality
says that it returns the continuous closest point of \(\mathcal V_B(\mathbb
R)\). This is still not discrete CVP unless a separate lattice is supplied.

\subsection{Explicit Harmonic-Lattice Bridge}

The discrete bridge requires both kinds of data: a real harmonic subspace and
an integer lattice contained in it. Let \(W\subseteq\mathbb R^m\) be the
operative harmonic subspace, let \(P_W\) be its orthogonal projection, and let
\(\Lambda_B=B\mathbb Z^n\subseteq W\) be an explicitly presented lattice. Set
\[
  \Delta_W=I-P_W.
\]
The inclusion \(\Lambda_B\subseteq W\) does not identify the discrete lattice
with the real subspace.

\begin{proposition}[Discrete CVP through harmonic projection]\label{prop:hodge-projection-is-cvp}
  For \(t\in\mathbb R^m\), put \(t_W=P_Wt\). Then:
  \begin{enumerate}[label=(\roman*)]
    \item \(t_W\) and every \(\lambda\in\Lambda_B\) are
      \(\Delta_W\)-harmonic;
    \item for every \(\lambda\in\Lambda_B\),
      \[
        \|t-\lambda\|^2
        =
        \|t-t_W\|^2+\|t_W-\lambda\|^2;
      \]
    \item a lattice vector \(\lambda\) is closest to \(t\) in
      \(\Lambda_B\) if and only if it is closest to \(t_W\) in
      \(\Lambda_B\).
  \end{enumerate}
\end{proposition}

\begin{proof}
  Since \(t_W\in W\), every \(\lambda\in\Lambda_B\) also lies in \(W\), and
  \(\ker(\Delta_W)=W\), the harmonicity claims follow. For the distance
  identity, write
  \[
    t-\lambda=(t-t_W)+(t_W-\lambda).
  \]
  Orthogonal projection gives \(t-t_W\in W^\perp\), while
  \(t_W-\lambda\in W\). The two summands are orthogonal, so the Pythagorean
  identity gives (ii). The term \(\|t-t_W\|^2\) is independent of
  \(\lambda\); consequently, for any \(\lambda,\mu\in\Lambda_B\),
  \[
    \|t-\lambda\|^2\le \|t-\mu\|^2
    \quad\Longleftrightarrow\quad
    \|t_W-\lambda\|^2\le \|t_W-\mu\|^2.
  \]
  This is exactly (iii).
\end{proof}

The continuous projection \(t_W=P_Wt\) is the closest harmonic vector to
\(t\), but it need not lie in \(\Lambda_B\). Thus projection transports the
target to the harmonic subspace while preserving the discrete minimizers; it
does not perform the remaining lattice search.

\subsection{Oracle Insufficiency}

Fix a typed HCP semantic realization as in \cref{def:typed-hcp-realization}.
The solver contracts below range over its admitted representatives.

\begin{definition}[Public-data Hodge solver contract]\label{def:public-hodge-solver-contract}
  Let
  \[
    S:\mathcal O_{\mathrm{CVP}}\times
      (P\times\mathbb R^n)\longrightarrow\mathbb R^n,
  \]
  where \(\mathcal O_{\mathrm{CVP}}\) is the class of explicit-lattice CVP
  oracles. For a realization \(x\mapsto(M_x,\Delta_x)\) and
  \(O\in\mathcal O_{\mathrm{CVP}}\), define
  \[
  \begin{aligned}
    \mathsf H_O(S)
    &\Longleftrightarrow
    \forall x:A,\quad
    \forall b\in\mathbb R^n,\quad
    \Delta_x\bigl(S(O,(\pi(x),b))\bigr)=0,\\
    \mathsf F_O(S)
    &\Longleftrightarrow
    \forall x:A,\quad
    \forall b\in\mathbb R^n,\quad
    \Delta_xb=0\Rightarrow S(O,(\pi(x),b))=b.
  \end{aligned}
  \]
  The public-data-only contract is
  \[
    \exists S,\quad
    \forall O,\quad
    \mathsf{CVP}_{\mathrm{corr}}(O)\Rightarrow
    \bigl(
      \mathsf H_O(S)
      \wedge
      \mathsf F_O(S)
    \bigr).
  \]
\end{definition}

\begin{proposition}[Typed HCP ambiguity and CVP oracle insufficiency]\label{prop:cvp-insufficient}
  Let \(x,y\in A\) be representatives in a typed HCP semantic realization.
  Suppose that \(M_x\) and \(M_y\) are homotopy equivalent and that some
  \(b\in\mathbb R^n\) satisfies
  \[
    \Delta_xb=0,
    \qquad
    \Delta_yb\ne0.
  \]
  If at least one correct explicit-lattice CVP oracle exists, then
  \[
    \neg\exists S,\quad
    \forall O,\quad
    \mathsf{CVP}_{\mathrm{corr}}(O)\Rightarrow
    \bigl(
      \mathsf H_O(S)
      \wedge
      \mathsf F_O(S)
    \bigr).
  \]
  Hence a correct oracle for explicit known-lattice CVP does not by itself
  determine the operative Hodge--Laplacian kernel.
\end{proposition}

\begin{proof}
  By \cref{prop:helfer-hcp-bridge},
  \(\pi(x)=\pi(y)\). Suppose
  \(\exists O,\mathsf{CVP}_{\mathrm{corr}}(O)\), and choose such an \(O\).
  Suppose, for contradiction, that \(S\) satisfies the displayed
  public-data-only contract. Then
  \[
    \mathsf H_O(S)\wedge\mathsf F_O(S).
  \]
  Since \(\Delta_xb=0\), exactness at \(x\) gives
  \[
    S(O,(\pi(x),b))=b.
  \]
  Soundness at \(y\) gives
  \[
    \Delta_yS(O,(\pi(y),b))=0.
  \]
  Replacing \(\pi(y)\) by \(\pi(x)\) and using the exactness equality yields
  \(\Delta_yb=0\), contradicting the separation hypothesis.
\end{proof}

For every dimension at least two, consider the elementary matrix witness \(L=0\),
\(L'=\operatorname{diag}(1,0,\ldots,0)\), and \(b=e_1\) under the constant
projection to a one-point public type. This is a sanity example showing that
the abstract ambiguity predicates are inhabited. It is not used as evidence
that an arbitrary HCP realization admits the required semantic and Hodge
extension witness. That witness remains an explicit hypothesis of
\cref{prop:cvp-insufficient}; the following family constructs one for its
stated semantic fragment.

\subsection{A Constructed Integral Cellular Family}

The preceding implication has a concrete instantiation in a stated
homological fragment of the typed Helfer interface. Take one vertex, one edge
and one face with incidence numbers \(\partial_1=0\), \(\partial_2=1\),
and stalks \(0,\mathbb R^2,\mathbb R\), respectively. This is the algebraic
finite signed-incidence model of \cref{sec:primitives}. Two cellular
sheaves in that model have
\[
  \delta_{0,0}=\delta_{0,1}=0,
  \qquad
  \delta_{1,0}(x_0,x_1)=x_0,
  \qquad
  \delta_{1,1}(x_0,x_1)=x_1.
\]
Their true degree-one Hodge operators are
\[
  \Delta_0=\operatorname{diag}(1,0),
  \qquad
  \Delta_1=\operatorname{diag}(0,1).
\]
Coordinate exchange is a degreewise chain isometry. Both operators have
positive spectral gap \(1\). Standard integral stalk lattices give harmonic
lattices \(\mathbb Z e_1\) and \(\mathbb Z e_0\), each with an explicit
full-span presentation.

Interpret these degree-zero-through-two complexes in the bounded-below
real cochain Quillen model category, with the incidence-defined coboundaries
displayed above. The closed tests are indexed by integers and ask whether
cohomology in that degree is nonzero. The required cofibrancy and fibrancy
hold for these real complexes. The model's weak equivalences are
quasi-isomorphisms, which preserve the tests and give the required
closed-test invariance under the declared typed Helfer equivalences.

The coordinate-exchange chain isometry gives the same homological theory
to both presentations. On this admitted family, equality of the public
theories is equivalent to typed homotopy equivalence.

\Needspace{10\baselineskip}
\begin{proposition}[One observation, two nonzero integral harmonic targets]
  \label{prop:cellular-axis-ambiguity}
  Let \(p\) be the common closed theory of these two presentations in the
  stated fragment, and let \(b=e_0+e_1\). There are structured Hodge samples
  with that same observation \((p,b)\), harmonic messages \(e_1\) and \(e_0\),
  and zero exact components. Their operative targets are distinct and nonzero.
  For each presentation's explicit harmonic lattice, its message is also the
  unique closest lattice vector to \(b\).
\end{proposition}

\begin{proof}
  In the first presentation use harmonic message \(e_1\), exact opening zero
  and coexact opening \(1\), producing noise \(e_0\). In the second use
  message \(e_0\), exact opening zero and coexact opening \(1\), producing
  noise \(e_1\). Both sums equal \(b\), and both messages belong to their
  standard integral harmonic lattices. The actual Hodge projections recover
  those messages. Orthogonality gives the Pythagorean identity from
  \cref{prop:hodge-projection-is-cvp}; any other lattice point has strictly
  positive squared distance from the recovered message. Thus each message
  is the unique CVP answer in its presented lattice.
\end{proof}

For the same realization, the fixed coordinate \(e_1\) is harmonic for
\(\Delta_0\) and not for \(\Delta_1\). This constructs the separation witness
used in \cref{prop:cvp-insufficient}. Equality of the public theory is
derived from the
actual semantic equivalence, rather than imposed by a one-point classifier.
The matrix-level public projection is extended from the two admitted
operators; no semantic realization is asserted for other matrices.

The displayed pair uses coefficient \(1\) in each presentation's own
harmonic-lattice basis. Its distinct ambient targets therefore do not alone
assert distinct encoded integer plaintexts. A second constructed pair makes
the encoding explicit: take \(b=e_0+2e_1\), harmonic messages \(2e_1\) and
\(e_0\), and complementary coexact noise \(e_0\) and \(2e_1\). Their integer
coefficients in the same declared lattice presentations are \(2\) and \(1\).
Both messages are nonzero, both exact components remain zero, and the public
observation is again shared. No function of that observation can return the
correct coefficient for both presentations. Each message remains the unique
CVP answer in its actual lattice.

\Needspace{8\baselineskip}
\begin{corollary}[Strict separation for hidden lattice recovery]
  \label{cor:lattice-member-separation}
  Fix a correct CVP oracle for the presented rank-one lattices of the
  constructed family. When the operative lattice is supplied, the oracle
  recovers its unique closest vector, equal to the operative harmonic target.
  From the shared public class and cochain alone, no solver with that same
  fixed oracle can recover the operative target for both hidden members.
  Every larger compatible family retaining those two members inherits this
  exact-recovery obstruction.
\end{corollary}

\begin{proof}
  For each supplied lattice, oracle correctness and
  \cref{prop:cellular-axis-ambiguity} identify its answer with the harmonic
  message. The two hidden members have the same public observation and
  distinct operative targets. A solver receives the same observation and
  the same oracle in either case, so its single output cannot equal both
  targets. The argument allows arbitrary fixed auxiliary information,
  including all candidate CVP answers, and imposes no running-time bound.
  Any compatible family containing the pair retains this contradiction.
\end{proof}

Solving the candidate CVP instances does not select the operative HCP answer.
